\documentclass
  [hyperref={colorlinks = true,linkcolor = blue, 
             citecolor = blue, urlcolor = blue}
  ]{beamer}

\setbeamertemplate{navigation symbols}{}
\setbeamertemplate{footline}[frame number]
 
\usepackage[utf8]{inputenc}

\usepackage[newfloat]{minted}

\usepackage{pgf}
\usepackage{tikz}
\usepackage{upquote}
\usepackage{natbib}
\usepackage[export]{adjustbox}
\usepackage{amsmath}
\usepackage{amssymb}
\usepackage{trfrac}

\usetikzlibrary{arrows,automata,fit, shapes.geometric}

\bibliographystyle{abbrvnat}

\newenvironment{code}{\captionsetup{type=listing}}{}
\SetupFloatingEnvironment{listing}{name=Listing}

\setbeamertemplate{items}[square]

\renewcommand{\thempfootnote}{\arabic{mpfootnote}}

\newcommand{\Set}{

}

\title{Why is so much of FP about Types instead of Functions?}
\author{Donovan Crichton}
  
\date{October 2023}

\begin{document}
 
\frame{\titlepage}

\begin{frame}[fragile]
  \frametitle{Preliminaries}
  \begin{itemize}
  \item Slides and Examples available at:
    \href{https://github.com/donovancrichton/Talks}
         {https://github.com/donovancrichton/Talks}
  \item This talk: BFPG/WhyIsFPAboutTypes
  \end{itemize}
\end{frame}

\begin{frame}[fragile]
\frametitle{About Me}
\begin{minipage}[t]{\linewidth}
\noindent
  \begin{minipage}{0.5\linewidth}
    \begin{figure}[h]
       \includegraphics[width=\linewidth]
         {Images/ANU.png}
    \end{figure}
  \end{minipage}
  \hfill
  \begin{minipage}{0.45\linewidth}
    \begin{itemize}
      \item 
        \href{https://cecc.anu.edu.au/people/donovan-crichton}
             {PhD Candidate}
      \item 
        \href{https://comp.anu.edu.au/research/clusters/computing-foundations/}
             {Computing Foundations}
      \item 
        \href{https://comp.anu.edu.au/}
             {School of Computing}
    \end{itemize}
  \end{minipage}
\end{minipage}
\hrule
  \begin{minipage}{0.5\linewidth}
    \begin{figure}[h]
      \includegraphics[scale=0.5,width=\linewidth]
        {Images/GriffithLogo.png}
    \end{figure}
  \end{minipage}
  \hfill
  \begin{minipage}{0.45\linewidth}
  \begin{itemize}
    \item Visiting Scholar
    \item Trusted Systems Lab
    \item \href{https://www.griffith.edu.au/institute-integrated-intelligent-systems}{IIIS}
  \end{itemize}
  \end{minipage}
\hrule
  \begin{minipage}{0.5\linewidth}
    \begin{figure}[h]
      \includegraphics[width=\linewidth]
        {Images/asd-logo.png}
    \end{figure}
  \end{minipage}
  \hfill
  \begin{minipage}{0.45\linewidth}
  \begin{itemize}
    \item ASD 
      \href{https://www.asd.gov.au/about/asd-anu-co-lab}{Co-Lab} 
      Scholar
  \end{itemize}
  \end{minipage}
\end{frame}

\begin{frame}[fragile]
\frametitle{What is a Function?}
\begin{minipage}{0.8\linewidth}
  \begin{figure}[h]
  \includegraphics
    [scale=0.3, width=0.6\linewidth]
    {Images/Function_machine.png}
  \caption[A "black-box" depiction of a function.]{A
  "black-box" depiction of a function\protect\footnote[1]{Source \protect\url{https://en.wikipedia.org/wiki/Function_(mathematics)\#/media/File:Function_machine2.svg}}}
  \end{figure}
\end{minipage}
\end{frame}

\begin{frame}[fragile]
\frametitle{What can we do with a function?}
  \begin{itemize}[<+->]
    \item Give it its argument.
    \item Look at its result.
    \item ...
    \item ...
    \item Maybe we need to think about this some more...
  \end{itemize}
\end{frame}

\begin{frame}[fragile]
\frametitle{What is a function, Georg Cantor?}
\begin{minipage}{0.8\linewidth}
  \begin{figure}[h]
  \includegraphics
    [scale=0.1, width=0.5\linewidth]
    {Images/Georg_Cantor.jpg}
  \caption[Georg Cantor (1845 - 1918)]{Georg Cantor (1845 -
  1918)\protect\footnote[1]{Source
  \protect\url{https://en.wikipedia.org/wiki/Georg_Cantor\#/media/File:Georg_Cantor_(Portr\%C3\%A4t).jpg}}}
  \end{figure}
\end{minipage}
\end{frame}

\begin{frame}[fragile]
\frametitle{Sets}
  \begin{itemize}[<+->]
    \item Sets are concerened with collections of discrete
          mathematical objects. \\
          e.g $A = \{1, 2, 3, ...\}$.
    \item The order of elements does not matter.
    \item Each element my appear only once.
    \item Given an set $X$ and an index $i$ we may 
          chose $x_i \in X$. Think of this like a function
          $f : (Set(X), \mathbb{N}) \rightarrow X$. \\
          e.g $f(A,2) = 2$.
    \item See \citep{halmos1960naive} for more!
  \end{itemize}
\end{frame}

\begin{frame}[fragile]
\frametitle{Relations}


\end{frame}

\begin{frame}[fragile]
\frametitle{Functions}
\end{frame}

\begin{frame}[fragile]
\frametitle{What can we do with set-theoretic functions?}
\end{frame}

\begin{frame}[fragile]
\frametitle{What can we do with functions, Alonzo Church?}
\begin{minipage}{0.8\linewidth}
  \begin{figure}[h]
  \includegraphics
    [scale=0.3, width=0.5\linewidth]
    {Images/Alonzo_Church.jpg}
  \caption[Alonzo Church (1903 - 1995)]{Alonzo Church (1903 -
  1995)\protect\footnote[1]{Source \protect\url{https://en.wikipedia.org/wiki/Alonzo_Church\#/media/File:Alonzo_Church.jpg}}}
  \end{figure}
\end{minipage}
\end{frame}

\begin{frame}[fragile]
\frametitle{The Untyped Lambda Calculus}  
  \setbeamercolor{block title}{fg=white,bg=purple!75!black}
  \setbeamercolor{block body}{fg=black,bg=white!80!gray}
  \begin{block}{The Grammar for LC.}
  \begin{align*}
    M, N ::= & \; x, y, z, ...  &\text{Variables.} \\
             &| \; \lambda x.N  &\text{Abstraction.} \\
             &| \; M \; N       &\text{Application}
  \end{align*}  
  \end{block}
  \setbeamercolor{block title}{fg=white,bg=purple!75!black}
  \setbeamercolor{block body}{fg=black,bg=white!80!gray}
  \begin{block}{An Idris Example (some liberties with types) 
  \\ $\lambda x.\lnot x \; \text{True}$}
  \inputminted{idris}{../Code/src/Lambda.idr}
\end{block}
\end{frame}

\begin{frame}[fragile]
\frametitle{What can we do with lambda calculi functions?}
\end{frame}

\begin{frame}[fragile]
\frametitle{Functions are boring!}
  \setbeamercolor{block title}{fg=black,bg=green!95!black}
  \setbeamercolor{block body}{fg=white,bg=green!25!black}
  \begin{block}{Can we apply any function to any argument?}
  Clearly we cannot, saying $(3 < 7) + 12$ does not make
  sense, nor does $ (-4) \lor 7$. 
  \\ Operations are often defined as being closed under 
  a particular set. \\
  Closures are really speaking about the type of operations.
  \end{block}
  \setbeamercolor{block title}{fg=black,bg=green!95!black}
  \setbeamercolor{block body}{fg=white,bg=green!25!black}
  \begin{block}{Are there other ways we can catgorise
    functions besides the type of their input and output?}
  We can also classify functions based on properties of their
  behaviour. Sometimes we don't care what specific types
  a function operates on, so long as there is some valid
  operation defined for those types.
  \end{block}
\end{frame}

\begin{frame}[fragile]
\frametitle{Diagrams of types.}
\end{frame}

\begin{frame}[fragile]
\frametitle{Diagrams of behavioural properties.}
\end{frame}

\begin{frame}[fragile]
\frametitle{To functional programming}
\setbeamercolor{block title}{fg=white,bg=orange!75!black}
\setbeamercolor{block body}{fg=white,bg=orange!25!black}
\begin{block}{Question for the pub!}
Why are functions the default 'building-blocks' in most
programming languages? Why not relations?
\end{block}

\setbeamercolor{block title}{fg=black,bg=green!95!black}
\setbeamercolor{block body}{fg=white,bg=green!25!black}
\begin{block}{Is everything really a function? (Really?)}
The argument goes that a pure functional language is just
an elaboration of some variant of a typed lambda calculus.
Let's briefly investigate this.
\end{block}
\end{frame}

\begin{frame}[fragile]
\frametitle{Values as functions.}
  \setbeamercolor{block title}{fg=white,bg=purple!75!black}
  \setbeamercolor{block body}{fg=black,bg=white!80!gray}
  \begin{block}{Are these functions the same?}<1->
  \inputminted{idris}{../Code/src/Values.idr}
  \end{block}

  \setbeamercolor{block title}{fg=black,bg=green!95!black}
  \setbeamercolor{block body}{fg=white,bg=green!25!black}
  \begin{block}{The humble unit type.}<2->
  The unit type is special.
  Denoted \mintinline{idris}{()} or $\top$, the unit type is
  defined as having only one single element, also
  \mintinline{idris}{()}, or sometimes $\ast$.
  The unit element can \textit{always} be constructed.
  This implies that \mintinline{idris}{seven} is the same as
  \mintinline{idris}{seven'} as we can always construct the 
  argument to \mintinline{idris}{seven'}.
  \end{block}
\end{frame}

\begin{frame}[fragile]
\frametitle{Data types as functions.}
  \setbeamercolor{block title}{fg=black,bg=green!95!black}
  \setbeamercolor{block body}{fg=black,bg=white!80!gray}
  \begin{block}{Can we do the same with data types?}<1->
    We can see this with church encoding, and have a
    small example below, interested readers should see
    \cite[p. 58-68]{pierce2002types}.
  \end{block}

  \setbeamercolor{block title}{fg=white,bg=purple!75!black}
  \setbeamercolor{block body}{fg=black,bg=white!80!gray}
  \begin{block}{Church encoding of the Boolean type}<2->
    \inputminted{idris}{../Code/src/ChurchBools.idr}
  \end{block}
\end{frame}

\begin{frame}[fragile]
\frametitle{Algerbraic Data Types}
  \setbeamercolor{block title}{fg=black,bg=green!95!black}
  \setbeamercolor{block body}{fg=white,bg=green!25!black}
  \begin{block}{Sum Types}
  The sum type is analogous to disjoint union in set theory.
  $A + B$ can be construct from either $x \in A$ or $y \in B$
  provided we tag the origin set.
  \end{block}

  \setbeamercolor{block title}{fg=black,bg=green!95!black}
  \setbeamercolor{block body}{fg=white,bg=green!25!black}
  \begin{block}{Product Types}
  
  \end{block}
\end{frame}

\begin{frame}[fragile]
\frametitle{Conclusion}

  \setbeamercolor{block title}{fg=black,bg=green!95!black}
  \setbeamercolor{block body}{fg=white,bg=green!25!black}
  \begin{block}{}
  \end{block}
\end{frame}
% question for the pub. Tension between 

% Introduce Georg Cantor for sets
%  - Arguably introduced in 1978
%  - No Bijection between Natural Numbers and real numbers?
%  - Both sets are infinite but there are more of one.

% Introduce Alonzo Church for lambda calculus
% Find out who introduced generics
% Find out who introduced parametric and ad-hoc polymorphism
% Introduce those two
% 

%Question for the pub: Why are 'functions' the basic building
%block in programming languages instead of relations?
%arguably relations are written just as often and are
% strict superset of functions!

%Question: Is everything really a function? Show 7 () :: Int
%instead of 7 :: Int


% Functions are boring!
% All we can do with them is apply them to their arguments.
% The second we start asking if we can apply functions to the
% result of other functions we composition.
% When we get composition, we start asking well does it make
% sense to compose two functions.
% Now we get types.
% Types cause tension between composability and safety.
% Even(N) does not compose with N but a domain taking an
% N for an input when we care about even N forces us to
% consider undefined behaviour.
% Functional Programming is about 'reducing unintended side
% affects"
% Undefined behaviour is an unintended side effect!



%in general I want to say that in programming

% Cover: zSets, Relations, Functions, Partial and Total
% Functions

% Cover: Lambda Calculus

% Cover: 

\begin{frame}[fragile]
\frametitle{References}
\bibliography{references}{}
\end{frame}




\end{document}

